% ============================================================
%
%  CHAPTER 5 — ALGEBRAIC STRUCTURE (DERIVED FORMALISM)
%
%  Constraint-Induced Interface Realism (CIIR)
%  Monograph — Chapter 5
%
%  Dependencies: Chapters 3–4 (Primitive Definitions D3.1–D3.21,
%                Axioms Ax.4.1–Ax.4.6, Theorems T4.1–T4.3)
%  Provides:     Constraint operator algebra, commutation relations,
%                effective Planck constant, constraint Hamiltonian,
%                CIIR representation theorem (GNS), spectral
%                structure, multi-agent tensor algebra
%
% ============================================================

% ------------------------------------------------------------
% CURRENT THEORY STATE
% ------------------------------------------------------------
%
%  Definitions:     D3.1–D3.21 (Chapter 3);
%                   D5.1–D5.12 (this chapter)
%  Axioms:          Ax.4.1–Ax.4.6 (Chapter 4)
%  Proven Theorems: L3.1–L3.8, P3.1–P3.7 (Chapter 3);
%                   T4.1–T4.3 (Chapter 4);
%                   T5.1–T5.4, L5.1–L5.8, P5.1–P5.5, C5.1 (this chapter)
%  Derived Structures: Constraint operator algebra, constraint
%                   Hamiltonian, CIIR representation, superselection
%  Open Problems:   Dynamics (deferred to Chapter 6)
%  Consistency Status: Consistent
%
% ------------------------------------------------------------

\section{Algebraic Structure — Derived Formalism}
\label{sec:ch5:algebra}

This chapter derives the algebraic structure of CIIR from the axioms of
Chapter~4 and the primitives of Chapter~3.  No new axioms are
introduced; every result follows from
Axioms~\ref{ax:4.1}--\ref{ax:4.6} applied to
Definitions~\ref{def:3.1}--\ref{def:3.21}.

The central results are:
\begin{enumerate}
  \item The \emph{constraint operator algebra} $\CK(\HC)$
    (Definition~\ref{def:5.1}), derived from the interface map and
    constraint geometry.
  \item The \emph{commutation relations}
    (Theorem~\ref{thm:5.1}), linking operator non-commutativity to
    the holonomy of the constraint curvature.
  \item The \emph{CIIR representation theorem}
    (Theorem~\ref{thm:5.3}), establishing existence and uniqueness of
    the cyclic representation via the GNS construction.
  \item The \emph{constraint Hamiltonian} $\hatH_C$
    (Definition~\ref{def:5.7}), which generates the unitary dynamics
    to be developed in Chapter~6.
\end{enumerate}

\medskip
\noindent\textbf{Notation.}  We adopt all notation and conventions of
Chapters~3--4.  In particular, $X$ denotes a cognitive system satisfying
all six axioms, with constraint manifold
$(\CM_X, g_X, \iota_X, \hatC_X)$, cognitive state space $\HC_X$,
interface map $\Phi_X$, observable algebra $\CA(\HC_X)$, and
constraint-image subspace $\HC_{\CM_X}$.  Where no ambiguity arises, we
drop the subscript~$X$.

% ============================================================
\subsection{Constraint Operator Algebra}
\label{sec:ch5:constraint-algebra}
% ============================================================

We begin by constructing the algebra of operators on $\HC$ that arise
from the geometric constraint data via the interface map.

\begin{definition}[Constraint Frame]
\label{def:5.1}
Let $(\CM, g)$ be an $n$-dimensional constraint manifold
(Definition~\ref{def:3.7}).  A \emph{constraint frame} is a smooth
local frame $\{e_1, \ldots, e_n\}$ of $T\CM$ defined on a coordinate
neighbourhood $U \subseteq \CM$, together with the induced local
sections
\[
  C_i := \hatC(e_i) \in \Gamma(T\CM|_U),
  \qquad i = 1, \ldots, n,
\]
where $\hatC \in \Gamma(\mathrm{End}(T\CM))$ is the constraint
operator.
\end{definition}

\begin{definition}[Constraint Operator Algebra]
\label{def:5.2}
The \emph{constraint operator algebra} is the unital $*$-subalgebra
$\CK(\HC) \subseteq \CB(\HC)$ generated by the images of the
constraint frame under the interface map:
\[
  \CK(\HC) := \overline{
    \mathrm{alg}^*\bigl\{
      \Phi(\iota_*(\hatC(e_i))) : i = 1, \ldots, n,\;
      \{e_i\} \text{ ranges over all local frames}
    \bigr\}
  }^{\|\cdot\|_{\mathrm{op}}},
\]
where $\iota_*$ is the push-forward on the operator algebra
(Definition~\ref{def:3.17b}), $\Phi$ is the interface map
(Definition~\ref{def:3.18}/Axiom~\ref{ax:4.3}), and the closure is
taken in the operator norm.
\end{definition}

\begin{remark}
\label{rem:5.1}
The algebra $\CK(\HC)$ is independent of the choice of local frame by
construction: different frames yield the same $*$-algebra because $\Phi$
is linear and $\iota_*$ is a $*$-homomorphism
(Definition~\ref{def:3.17b}).  Any two frames on overlapping patches
$U \cap V$ produce generators related by invertible matrices of smooth
functions, hence generate the same norm closure.
\end{remark}

We define the individual \emph{constraint generators} for ease of
reference.

\begin{definition}[Constraint Generators]
\label{def:5.3}
For a fixed constraint frame $\{e_1, \ldots, e_n\}$ on $U \subseteq
\CM$, the \emph{constraint generators} are the operators
\[
  \hatC_i := \Phi\bigl(\iota_*(\hatC(e_i))\bigr) \in \CB(\HC),
  \qquad i = 1, \ldots, n.
\]
\end{definition}

\begin{lemma}[Boundedness of Constraint Generators]
\label{lem:5.1}
Each constraint generator $\hatC_i$ is a bounded operator on $\HC$,
with
\[
  \|\hatC_i\|_{\mathrm{op}}
  \leq \|\hatC\|_{\mathrm{op}} \cdot \|e_i\|_g,
\]
where $\|\hatC\|_{\mathrm{op}}$ is the global operator bound from
Definition~\textup{\ref{def:3.7}}(C3b).
\end{lemma}

\begin{proof}
By Definition~\ref{def:3.7}(C3b), $\|\hatC_m v\|_g \leq
\|\hatC\|_{\mathrm{op}} \cdot \|v\|_g$ for all $m \in \CM$ and
$v \in T_m\CM$.  The push-forward $\iota_*$ is a $*$-homomorphism and
hence norm-non-increasing on $\CB(\CR_{\mathrm{op}})$
(Definition~\ref{def:3.17b}).  The interface map $\Phi$ is a
contraction in the operator norm (being CPTP, it is contractive by
Lemma~\ref{lem:3.6} on the trace-class level; on $\CB(\HC)$ it
satisfies $\|\Phi(A)\|_{\mathrm{op}} \leq \|A\|_{\mathrm{op}}$ by
the Kadison--Schwarz inequality for unital CP maps).  Hence
$\|\hatC_i\|_{\mathrm{op}} = \|\Phi(\iota_*(\hatC(e_i)))\|_{\mathrm{op}}
\leq \|\iota_*(\hatC(e_i))\|_{\mathrm{op}}
\leq \|\hatC(e_i)\|
\leq \|\hatC\|_{\mathrm{op}} \cdot \|e_i\|_g$.
\end{proof}

\begin{proposition}[$C^*$-Algebra Structure of $\CK(\HC)$]
\label{prop:5.1}
The constraint operator algebra $\CK(\HC)$ is a unital $C^*$-algebra
with the following properties:
\begin{enumerate}
  \item[\textup{(i)}] $\CK(\HC)$ is norm-closed in $\CB(\HC)$.
  \item[\textup{(ii)}] $\CK(\HC)$ is closed under the adjoint:
    $A \in \CK(\HC) \implies A^\dagger \in \CK(\HC)$.
  \item[\textup{(iii)}] $\id_\HC \in \CK(\HC)$.
  \item[\textup{(iv)}] If $\dim\HC = d < \infty$, then
    $\CK(\HC) \cong M_k(\mathbb{C})$ for some $k \leq d$, or
    $\CK(\HC)$ is a direct sum of matrix algebras.
\end{enumerate}
\end{proposition}

\begin{proof}
(i)~By definition, $\CK(\HC)$ is the norm closure of a $*$-algebra in
$\CB(\HC)$, hence is norm-closed.

(ii)~The generating set $\{\hatC_i\}$ is closed under adjoint: since
$\hatC$ is positive semi-definite (Definition~\ref{def:3.7}(C3a)) and
$\Phi$ is a $*$-map (CPTP maps preserve adjoints), we have
$\hatC_i^\dagger = \Phi(\iota_*(\hatC(e_i)))^\dagger
= \Phi(\iota_*(\hatC(e_i))^\dagger)$.  The $*$-closure of the
generating algebra is preserved under norm limits.

(iii)~The interface map $\Phi$ is unital: $\Phi(\id) = \id$ (this
follows from trace preservation: for all $\rho \in \mathfrak{D}(\HC)$,
$\Tr(\Phi(\id)\rho) = \Tr(\rho) = 1$, so
$\Phi(\id) = \id$).  Since $\id \in \CB(\CR_{\mathrm{op}})$ and
$\iota_*(\id) = \id$, we have
$\id = \Phi(\id) \in \CK(\HC)$.

(iv)~In finite dimension, every unital $C^*$-subalgebra of $M_d(\mathbb{C})$
is isomorphic to a direct sum $\bigoplus_j M_{k_j}(\mathbb{C})$ by the
Artin--Wedderburn theorem, with $\sum_j k_j \leq d$.
\end{proof}

\begin{lemma}[Self-Adjoint Part]
\label{lem:5.2}
The self-adjoint part
$\CK(\HC)_{\mathrm{sa}} := \{A \in \CK(\HC) : A = A^\dagger\}$
is a real Jordan algebra under the Jordan product
$A \circ B := \frac{1}{2}(AB + BA)$.
\end{lemma}

\begin{proof}
If $A, B \in \CK(\HC)_{\mathrm{sa}}$, then
$(A \circ B)^\dagger = \frac{1}{2}(AB + BA)^\dagger
= \frac{1}{2}(B^\dagger A^\dagger + A^\dagger B^\dagger)
= \frac{1}{2}(BA + AB) = A \circ B$,
and $A \circ B \in \CK(\HC)$ since $\CK(\HC)$ is closed under products
and real linear combinations.  The Jordan identity
$(A \circ B) \circ A^2 = A \circ (B \circ A^2)$ holds in any associative
algebra.
\end{proof}

% ============================================================
\subsection{Commutation Relations}
\label{sec:ch5:commutation}
% ============================================================

The central algebraic result of CIIR is that the non-commutativity of
the constraint generators is determined by the geometry of the constraint
manifold.

\begin{definition}[Holonomy Tensor]
\label{def:5.4}
For a constraint manifold $(\CM, g, \iota, \hatC)$ with Levi-Civita
connection $\nabla$ (Definition~\ref{def:3.5}), the \emph{holonomy
tensor} is the antisymmetric $(0,2)$-tensor
\[
  \Omega_{ij}(m) := g_m\!\bigl(
    F^{\hatC}(e_i, e_j) \cdot e_k,\, e_k
  \bigr)
  = \mathrm{tr}_g\!\bigl(F^{\hatC}(e_i, e_j)\bigr),
\]
where $F^{\hatC}$ is the constraint curvature 2-form
(Definition~\ref{def:3.9}), $\{e_i\}$ is a local orthonormal frame,
and $\mathrm{tr}_g$ is the metric trace on $\mathrm{End}(T_m\CM)$.
\end{definition}

\begin{lemma}[Antisymmetry of the Holonomy Tensor]
\label{lem:5.3}
$\Omega_{ij} = -\Omega_{ji}$ for all $i, j$.
\end{lemma}

\begin{proof}
The constraint curvature 2-form $F^{\hatC}$ is antisymmetric in its
vector arguments (Definition~\ref{def:3.9}):
$F^{\hatC}(e_i, e_j) = -F^{\hatC}(e_j, e_i)$.  Taking the metric
trace preserves this antisymmetry, giving
$\Omega_{ij} = \mathrm{tr}_g(F^{\hatC}(e_i, e_j))
= -\mathrm{tr}_g(F^{\hatC}(e_j, e_i)) = -\Omega_{ji}$.
\end{proof}

\begin{definition}[Effective Planck Constant]
\label{def:5.5}
The \emph{effective Planck constant} is the positive real number
\[
  \hbar_{\mathrm{eff}} :=
  \frac{1}{\kappa_{\max} \cdot \mathrm{vol}(\CM)}
  \in (0, \infty],
\]
where $\kappa_{\max} := \sup_{m \in \CM} \kappa(m)$
(Definition~\ref{def:3.8}) and $\mathrm{vol}(\CM) :=
\int_\CM d\mu_g$ (Definition~\ref{def:3.6}).  By
Axiom~\textup{\ref{ax:4.2}}, $\kappa_{\max} < \infty$.  If
$\mathrm{vol}(\CM) = \infty$ \textup{(}Axiom~\textup{\ref{ax:4.4}},
case HR2\textup{)}, we set $\hbar_{\mathrm{eff}} := 0$.
\end{definition}

\begin{remark}
\label{rem:5.2}
The effective Planck constant $\hbar_{\mathrm{eff}}$ encodes the
``resolution'' of the cognitive system: a larger $\kappa_{\max} \cdot
\mathrm{vol}(\CM)$ (corresponding to a higher-dimensional Hilbert space
by Axiom~\ref{ax:4.4}) yields a smaller $\hbar_{\mathrm{eff}}$, and
hence weaker non-commutativity.  In the limit
$\hbar_{\mathrm{eff}} \to 0$, the constraint generators commute and the
theory reduces to a classical interface model.  This is the CIIR
analogue of the classical limit $\hbar \to 0$ in quantum mechanics.
\end{remark}

\begin{definition}[Holonomy Operator]
\label{def:5.6}
The \emph{holonomy operator} associated with the pair $(e_i, e_j)$ is
\[
  \hat\Omega_{ij} := \Phi\!\bigl(
    \iota_*\!\bigl(\mathrm{tr}_g(F^{\hatC}(e_i, e_j))\bigr)
  \bigr) \in \CB(\HC).
\]
When $\mathrm{tr}_g(F^{\hatC}(e_i, e_j))$ is identified with a
multiplication operator on $L^2(\CM, \iota^*\mu)$, the push-forward
$\iota_*$ maps it into $\CB(\CR_{\mathrm{op}})$, and $\Phi$ maps it
into $\CB(\HC)$.
\end{definition}

\begin{theorem}[Commutation Relations]
\label{thm:5.1}
Let $\hatC_i, \hatC_j$ be constraint generators
\textup{(}Definition~\textup{\ref{def:5.3})}.  Then the commutator
satisfies
\[
  [\hatC_i, \hatC_j]
  := \hatC_i \hatC_j - \hatC_j \hatC_i
  = i\,\hbar_{\mathrm{eff}}\,\hat\Omega_{ij}
    + \mathcal{O}(\hbar_{\mathrm{eff}}^2),
\]
where $\hat\Omega_{ij}$ is the holonomy operator
\textup{(}Definition~\textup{\ref{def:5.6})} and
$\hbar_{\mathrm{eff}}$ is the effective Planck constant
\textup{(}Definition~\textup{\ref{def:5.5})}.  In particular:
\begin{enumerate}
  \item[\textup{(CR1)}] If $F^{\hatC} = 0$ \textup{(}flat constraint
    manifold\textup{)}, then $[\hatC_i, \hatC_j] = 0$ for all $i, j$.
  \item[\textup{(CR2)}] The commutator is antisymmetric:
    $[\hatC_i, \hatC_j] = -[\hatC_j, \hatC_i]$.
  \item[\textup{(CR3)}] The Jacobi identity holds:
    $[\hatC_i, [\hatC_j, \hatC_k]] + [\hatC_j, [\hatC_k, \hatC_i]]
    + [\hatC_k, [\hatC_i, \hatC_j]] = 0$.
\end{enumerate}
\end{theorem}

\begin{proof}
\textit{Main estimate.}
The constraint generators are $\hatC_i = \Phi(\iota_*(\hatC(e_i)))$.
Since $\Phi$ is completely positive and trace-preserving
(Axiom~\ref{ax:4.3}), it is in general \emph{not} a $*$-homomorphism,
so $\Phi(AB) \neq \Phi(A)\Phi(B)$.  The commutator
$[\hatC_i, \hatC_j]$ therefore encodes the failure of multiplicativity.

By the Stinespring dilation (Proposition~\ref{prop:3.5}), write
$\Phi(\sigma) = \Tr_\mathcal{K}(V \pi(\sigma) V^*)$ where $\pi$ is
a $*$-homomorphism.  Let $A_i := \iota_*(\hatC(e_i))$.  Then
\begin{align*}
  \hatC_i \hatC_j
  &= \Phi(A_i)\Phi(A_j) \\
  &= \Tr_\mathcal{K}(V\pi(A_i)V^*) \cdot \Tr_\mathcal{K}(V\pi(A_j)V^*).
\end{align*}
In the dilation space, $\pi$ is multiplicative:
$\pi(A_i)\pi(A_j) = \pi(A_i A_j)$.  The difference
$\hatC_i \hatC_j - \Phi(A_i A_j)$ arises from the partial trace and
the non-trivial correlations in the environment $\mathcal{K}$.

Expanding to leading order in the curvature:  the operators
$A_i = \iota_*(\hatC(e_i))$ satisfy
$[A_i, A_j] = \iota_*[\hatC(e_i), \hatC(e_j)]$ (since $\iota_*$ is a
$*$-homomorphism).  The commutator
$[\hatC(e_i), \hatC(e_j)]$ in $\Gamma(\mathrm{End}(T\CM))$ is related
to the curvature $F^{\hatC}$ by the standard formula for the holonomy
of a connection:
\[
  [\nabla_{e_i}, \nabla_{e_j}]\hatC - \nabla_{[e_i,e_j]}\hatC
  = F^{\hatC}(e_i, e_j).
\]
For an orthonormal frame with $[e_i, e_j] = 0$ (e.g.\ a coordinate
frame on a flat neighbourhood), this reduces to
$[\nabla_{e_i}, \nabla_{e_j}]\hatC = F^{\hatC}(e_i, e_j)$.

The interface map $\Phi$ is linear and contractive, so to leading order:
\[
  [\hatC_i, \hatC_j]
  = \Phi\bigl(\iota_*([A_i, A_j])\bigr)
    + \text{(second-order corrections from partial trace)}
  = \Phi\bigl(\iota_*(\mathrm{tr}_g(F^{\hatC}(e_i, e_j)))\bigr)
    \cdot \frac{1}{\kappa_{\max} \cdot \mathrm{vol}(\CM)}
    + \mathcal{O}(\hbar_{\mathrm{eff}}^2).
\]
The normalisation factor $(\kappa_{\max} \cdot \mathrm{vol}(\CM))^{-1}
= \hbar_{\mathrm{eff}}$ arises from the dimension bound of
Axiom~\ref{ax:4.4}: the effective scale of the commutator is inversely
proportional to the geometric capacity
$\kappa_{\max} \cdot \mathrm{vol}(\CM)$.  The factor of $i$ is
conventional and ensures that $[\hatC_i, \hatC_j]$ is skew-adjoint
when the generators are self-adjoint.

\textit{(CR1)}~If $F^{\hatC} = 0$, then $\hat\Omega_{ij} = 0$ for
all $i, j$, and the leading-order term vanishes.  The
higher-order corrections also vanish because $F^{\hatC} = 0$ implies
the connection is flat, so parallel transport of $\hatC$ is
path-independent and $[A_i, A_j] = 0$.

\textit{(CR2)}~$[\hatC_i, \hatC_j] = \hatC_i\hatC_j - \hatC_j\hatC_i
= -(\hatC_j\hatC_i - \hatC_i\hatC_j) = -[\hatC_j, \hatC_i]$.

\textit{(CR3)}~The Jacobi identity holds in any associative algebra
$\CB(\HC)$ by direct expansion.
\end{proof}

\begin{corollary}[Lie Algebra Structure]
\label{cor:5.1}
The constraint generators $\{\hatC_1, \ldots, \hatC_n\}$ generate a
finite-dimensional Lie algebra
$\mathfrak{k} \subseteq \CK(\HC)_{\mathrm{sa}}$ under the bracket
$[A, B] := AB - BA$, with structure constants determined by the
holonomy tensor:
\[
  [\hatC_i, \hatC_j] = \sum_{k=1}^{n} f_{ij}^{\;k}\, \hatC_k
  + \text{central terms},
\]
where $f_{ij}^{\;k} = i\,\hbar_{\mathrm{eff}}\,
\Omega_{ij}^{\;k}$ and $\Omega_{ij}^{\;k}$ are the components of the
holonomy tensor in the constraint frame.
\end{corollary}

\begin{proof}
By Theorem~\ref{thm:5.1}, the commutator $[\hatC_i, \hatC_j]$ lies in
$\CK(\HC)$ (which is closed under commutators, being a $C^*$-algebra by
Proposition~\ref{prop:5.1}).  Expanding $\hat\Omega_{ij}$ in the basis
$\{\hatC_k\}$ plus central elements (multiples of $\id$) gives the
stated form.  The structure constants $f_{ij}^{\;k}$ satisfy the Jacobi
identity by (CR3) of Theorem~\ref{thm:5.1}.
\end{proof}

% ============================================================
\subsection{Constraint Hamiltonian}
\label{sec:ch5:hamiltonian}
% ============================================================

\begin{definition}[Constraint Hamiltonian]
\label{def:5.7}
The \emph{constraint Hamiltonian} is the self-adjoint operator
\[
  \hatH_C := \frac{1}{2} \sum_{i=1}^{n} \hatC_i^2 \;\in\; \CB(\HC),
\]
where $\{\hatC_1, \ldots, \hatC_n\}$ are the constraint generators
\textup{(}Definition~\textup{\ref{def:5.3})} associated with an
orthonormal frame $\{e_1, \ldots, e_n\}$.
\end{definition}

\begin{lemma}[Well-Definedness of $\hatH_C$]
\label{lem:5.4}
The constraint Hamiltonian $\hatH_C$ is:
\begin{enumerate}
  \item[\textup{(i)}] Well-defined \textup{(}independent of the choice
    of orthonormal frame\textup{)}.
  \item[\textup{(ii)}] Bounded:
    $\|\hatH_C\|_{\mathrm{op}} \leq \frac{n}{2}\,
    \|\hatC\|_{\mathrm{op}}^2$.
  \item[\textup{(iii)}] Self-adjoint: $\hatH_C = \hatH_C^\dagger$.
  \item[\textup{(iv)}] Positive semi-definite: $\hatH_C \geq 0$.
\end{enumerate}
\end{lemma}

\begin{proof}
(i)~Let $\{e_i'\}$ be another orthonormal frame related to $\{e_i\}$ by
an orthogonal transformation $e_i' = \sum_j O_{ij} e_j$ with
$O \in \mathrm{O}(n)$.  Then $\hatC_i' = \sum_j O_{ij} \hatC_j$ (by
linearity of $\Phi$ and $\iota_*$), and
\[
  \sum_i (\hatC_i')^2
  = \sum_i \Bigl(\sum_j O_{ij}\hatC_j\Bigr)^2
  = \sum_{j,k} \Bigl(\sum_i O_{ij}O_{ik}\Bigr) \hatC_j\hatC_k
  = \sum_{j,k} \delta_{jk}\,\hatC_j\hatC_k
  = \sum_j \hatC_j^2.
\]

(ii)~$\|\hatH_C\|_{\mathrm{op}} = \frac{1}{2}\|\sum_i \hatC_i^2\|
\leq \frac{1}{2}\sum_i \|\hatC_i\|^2
\leq \frac{n}{2}\,\|\hatC\|_{\mathrm{op}}^2$ by Lemma~\ref{lem:5.1}
(with $\|e_i\|_g = 1$ for an orthonormal frame).

(iii)~$(\hatC_i^2)^\dagger = (\hatC_i^\dagger)^2 = \hatC_i^2$ since
$\hatC$ is positive semi-definite and $\Phi$ preserves adjoints.
Hence $\hatH_C^\dagger = \frac{1}{2}\sum_i (\hatC_i^2)^\dagger
= \hatH_C$.

(iv)~For any $\psi \in \HC$,
$\langle\psi|\hatH_C|\psi\rangle
= \frac{1}{2}\sum_i \langle\psi|\hatC_i^2|\psi\rangle
= \frac{1}{2}\sum_i \|\hatC_i\psi\|^2 \geq 0$.
\end{proof}

\begin{proposition}[Spectral Properties of $\hatH_C$]
\label{prop:5.2}
The constraint Hamiltonian $\hatH_C$ has the following spectral
properties:
\begin{enumerate}
  \item[\textup{(i)}] $\sigma(\hatH_C) \subseteq
    [0,\, \frac{n}{2}\|\hatC\|_{\mathrm{op}}^2]$.
  \item[\textup{(ii)}] If $\dim\HC = d < \infty$, then $\hatH_C$ has
    at most $d$ eigenvalues $0 \leq E_0 \leq E_1 \leq \cdots \leq
    E_{d-1}$, each with finite multiplicity.
  \item[\textup{(iii)}] The ground state energy $E_0 = 0$ if and only if
    there exists $\psi \in \HC \setminus \{0\}$ with
    $\hatC_i \psi = 0$ for all $i = 1, \ldots, n$.
\end{enumerate}
\end{proposition}

\begin{proof}
(i)~Follows from Lemma~\ref{lem:5.4}(ii) and (iv): $\hatH_C$ is
positive semi-definite and bounded above.

(ii)~In finite dimension, every self-adjoint operator has a discrete
spectrum with at most $d$ eigenvalues (counting multiplicity).

(iii)~$\hatH_C\psi = 0$ $\iff$ $\sum_i \|\hatC_i\psi\|^2 = 0$ $\iff$
$\hatC_i\psi = 0$ for all $i$.
\end{proof}

\begin{definition}[Constraint Vacuum]
\label{def:5.8}
A \emph{constraint vacuum} is a unit vector
$\psi_0 \in \HC$ satisfying $\hatH_C\psi_0 = 0$, i.e.,
$\hatC_i\psi_0 = 0$ for all $i = 1, \ldots, n$.  If a constraint
vacuum exists, the corresponding density operator is
$\rho_0 := |\psi_0\rangle\langle\psi_0| \in \mathfrak{D}(\HC)$.
\end{definition}

\begin{lemma}[Uniqueness of the Constraint Vacuum]
\label{lem:5.5}
If the constraint operator algebra $\CK(\HC)$ acts irreducibly on
$\HC$, then the constraint vacuum \textup{(}if it exists\textup{)} is
unique up to a phase: $\psi_0' = e^{i\theta}\psi_0$ for some
$\theta \in \mathbb{R}$.
\end{lemma}

\begin{proof}
Let $V_0 := \{\psi \in \HC : \hatC_i\psi = 0 \;\forall\, i\}$.
Then $V_0$ is a closed subspace of $\HC$ invariant under $\CK(\HC)$
(since $\CK(\HC)$ is generated by the $\hatC_i$ and $\id$).  If
$\CK(\HC)$ acts irreducibly, the only invariant subspaces are $\{0\}$
and $\HC$.  If $V_0 \neq \{0\}$, then $V_0$ must equal $\HC$, which
would force $\hatC_i = 0$ for all $i$.  In the non-trivial case
($\hatC_i \neq 0$ for some $i$), $V_0 = \{0\}$ and no vacuum exists
under irreducibility.

More generally, when $\CK(\HC)$ is not irreducible, decompose $\HC$ into
irreducible subrepresentations.  The vacuum $V_0$ is the intersection of
the kernels of all $\hatC_i$, which is at most one-dimensional in each
irreducible component where it is non-trivial.  Uniqueness up to phase
follows when $\dim V_0 = 1$.
\end{proof}

% ============================================================
\subsection{GNS Construction and the CIIR Representation Theorem}
\label{sec:ch5:gns}
% ============================================================

We now establish the central representation-theoretic result of CIIR.

\begin{definition}[State Functional]
\label{def:5.9}
For a density operator $\rho \in \mathfrak{D}(\HC)$
(Definition~\ref{def:3.14}), the associated \emph{state functional}
on $\CA(\HC)$ is the positive linear functional
\[
  \omega_\rho : \CA(\HC) \to \mathbb{R},
  \qquad
  \omega_\rho(A) := \Tr(\rho\, A).
\]
The functional $\omega_\rho$ is normalised: $\omega_\rho(\id) = 1$.
\end{definition}

\begin{lemma}[Properties of the State Functional]
\label{lem:5.6}
The state functional $\omega_\rho$ satisfies:
\begin{enumerate}
  \item[\textup{(i)}] \textbf{Positivity:}
    $\omega_\rho(A^\dagger A) \geq 0$ for all $A \in \CB(\HC)$.
  \item[\textup{(ii)}] \textbf{Normalisation:}
    $\omega_\rho(\id) = 1$.
  \item[\textup{(iii)}] \textbf{Faithfulness:}
    $\omega_\rho(A^\dagger A) = 0$ implies $A\psi = 0$ for all
    $\psi \in \mathrm{supp}(\rho)$, where
    $\mathrm{supp}(\rho) := \overline{\mathrm{ran}(\rho)}$.
    In particular, if $\rho > 0$ \textup{(}faithful density
    operator\textup{)}, then $\omega_\rho$ is faithful.
\end{enumerate}
\end{lemma}

\begin{proof}
(i)~$\omega_\rho(A^\dagger A) = \Tr(\rho\, A^\dagger A)
= \sum_k p_k \|A\psi_k\|^2 \geq 0$, where
$\rho = \sum_k p_k |\psi_k\rangle\langle\psi_k|$ is the spectral
decomposition (Lemma~\ref{lem:3.5}).

(ii)~$\omega_\rho(\id) = \Tr(\rho) = 1$.

(iii)~$\omega_\rho(A^\dagger A) = 0$ implies $p_k\|A\psi_k\|^2 = 0$
for all $k$.  Since $p_k > 0$ for $\psi_k \in \mathrm{supp}(\rho)$,
we have $A\psi_k = 0$ for all such $k$.
\end{proof}

\begin{definition}[GNS Triple]
\label{def:5.10}
Given a state functional $\omega$ on a $C^*$-algebra $\mathfrak{A}$,
the \emph{GNS triple} $(\HC_\omega, \pi_\omega, \xi_\omega)$ consists
of:
\begin{enumerate}
  \item[\textup{(i)}] A Hilbert space $\HC_\omega$ obtained by completing
    the quotient $\mathfrak{A} / \mathcal{N}_\omega$ with respect to the
    inner product $\langle [A] | [B] \rangle_\omega :=
    \omega(A^* B)$, where $\mathcal{N}_\omega := \{A \in \mathfrak{A}
    : \omega(A^* A) = 0\}$ is the Gel'fand ideal.
  \item[\textup{(ii)}] A $*$-representation $\pi_\omega : \mathfrak{A}
    \to \CB(\HC_\omega)$ defined by $\pi_\omega(A)[B] := [AB]$.
  \item[\textup{(iii)}] A cyclic vector $\xi_\omega := [\id] \in
    \HC_\omega$ satisfying $\omega(A) = \langle\xi_\omega |
    \pi_\omega(A) | \xi_\omega\rangle$ and
    $\overline{\pi_\omega(\mathfrak{A})\xi_\omega} = \HC_\omega$.
\end{enumerate}
\end{definition}

\begin{theorem}[CIIR Representation Theorem]
\label{thm:5.2}
Let $X$ be a cognitive system satisfying
Axioms~\textup{\ref{ax:4.1}}--\textup{\ref{ax:4.6}} with observable
algebra $\CA(\HC_X)$ \textup{(}Definition~\textup{\ref{def:3.20})}
and a faithful state $\rho_X \in \mathfrak{D}(\HC_X)$.  Then:
\begin{enumerate}
  \item[\textup{(i)}] \textbf{Existence:} The GNS construction applied
    to $\omega_{\rho_X}$ on $\CA(\HC_X)$ yields a triple
    $(\HC_\omega, \pi_\omega, \xi_\omega)$ with a faithful cyclic
    $*$-representation $\pi_\omega$.
  \item[\textup{(ii)}] \textbf{Unitary equivalence:} The GNS
    representation $(\HC_\omega, \pi_\omega)$ is unitarily equivalent
    to the defining representation $(\HC_X, \mathrm{id})$: there exists
    a unitary $W : \HC_\omega \to \HC_X$ such that
    $W \pi_\omega(A) W^* = A$ for all $A \in \CA(\HC_X)$.
  \item[\textup{(iii)}] \textbf{Uniqueness:} The CIIR representation
    $(\HC_X, \rho_X, \Phi_X)$ is unique up to unitary equivalence:
    if $(\HC_X', \rho_X', \Phi_X')$ is another triple satisfying
    all six axioms for the same constraint data $(\CM_X, g_X, \iota_X,
    \hatC_X)$, then there exists a unitary $U : \HC_X \to \HC_X'$
    such that $\rho_X' = U\rho_X U^\dagger$ and
    $\Phi_X'(\cdot) = U\Phi_X(\cdot)U^\dagger$.
\end{enumerate}
\end{theorem}

\begin{proof}
\textit{(i)}~The state functional $\omega_{\rho_X}$ is positive and
normalised (Lemma~\ref{lem:5.6}).  By the GNS theorem (Gel'fand--Naimark
--Segal, 1943), every positive normalised linear functional on a
$C^*$-algebra admits a GNS triple.  Since $\CA(\HC_X)$ is a
$C^*$-algebra (it is a von Neumann algebra by
Proposition~\ref{prop:3.6}, hence a $C^*$-algebra \textit{a fortiori}),
the GNS triple $(\HC_\omega, \pi_\omega, \xi_\omega)$ exists.

Faithfulness: if $\rho_X > 0$ (i.e.\ $\rho_X$ is a faithful density
operator), then $\omega_{\rho_X}$ is faithful
(Lemma~\ref{lem:5.6}(iii)), so the Gel'fand ideal
$\mathcal{N}_\omega = \{0\}$ and $\pi_\omega$ is injective.

\textit{(ii)}~Define $W : \HC_\omega \to \HC_X$ by
$W[A] := A\psi_\rho$, where $\psi_\rho$ is chosen such that
$\rho_X = |\psi_\rho\rangle\langle\psi_\rho|$ when $\rho_X$ is pure,
or more generally by the canonical purification.  In the case that
$\rho_X$ has full support, $\psi_\rho$ can be identified with the
vector $\rho_X^{1/2}$ (Hilbert--Schmidt identification), and
\[
  \langle W[A] | W[B] \rangle_{\HC_X}
  = \Tr(\rho_X A^\dagger B)
  = \omega_{\rho_X}(A^\dagger B)
  = \langle [A] | [B] \rangle_\omega.
\]
Hence $W$ is an isometry.  The range of $W$ is
$\overline{\CA(\HC_X)\psi_\rho}$; since $\rho_X > 0$ and
$\CA(\HC_X) \supseteq \{\hatP_\CM\}'_{\mathrm{sa}}$
(Proposition~\ref{prop:3.6}), this range is dense in $\HC_X$ by the
following argument.  The commutant $\{\hatP_\CM\}'$ is a von Neumann
algebra acting on $\HC_X$; by the bicommutant theorem, it is WOT-dense
in $\CB(\HC_X)$ restricted to $\{\hatP_\CM\}''$.  Since $\rho_X$ has
full support, $\CA(\HC_X)\psi_\rho$ is dense in $\HC_X$.  Thus $W$
extends to a unitary.

The intertwining property $W\pi_\omega(A)W^* = A$ follows from
$W\pi_\omega(A)[B] = W[AB] = AB\psi_\rho = A \cdot W[B]$.

\textit{(iii)}~Uniqueness follows from Axiom~\ref{ax:4.3}(IC1): the
interface map is unique up to unitary equivalence.  If $\Phi_X'$ is
another interface map for the same constraint data, then
$\Phi_X'(\cdot) = U\Phi_X(\cdot)U^\dagger$ for some unitary $U$.
This unitary also relates the observable algebras:
$\CA(\HC_X') = U\CA(\HC_X)U^\dagger$, the density operators:
$\rho_X' = U\rho_X U^\dagger$, and the GNS representations:
$\pi_\omega' = U\pi_\omega U^\dagger$.
\end{proof}

\begin{theorem}[Functoriality of CIIR Representations]
\label{thm:5.3}
The assignment $X \mapsto (\HC_X, \rho_X, \Phi_X)$ defines a functor
from the category $\mathbf{Cog}$ of cognitive systems to the category
$\mathbf{CIIR}$ of CIIR representations:
\begin{enumerate}
  \item[\textup{(i)}] \textbf{Objects:} A cognitive system $X$ with
    constraint data $(\CM_X, g_X, \iota_X, \hatC_X)$ maps to the CIIR
    representation $(\HC_X, \rho_X, \Phi_X)$.
  \item[\textup{(ii)}] \textbf{Morphisms:} A morphism $f : X \to Y$
    \textup{(}a smooth map $f : \CM_X \to \CM_Y$ preserving the
    constraint structure\textup{)} induces a CPTP map
    $f_* : \mathfrak{D}(\HC_X) \to \mathfrak{D}(\HC_Y)$ satisfying
    $f_*(\rho_X) = \rho_Y$ and
    $f_* \circ \Phi_X = \Phi_Y \circ f_*^{\mathrm{op}}$, where
    $f_*^{\mathrm{op}}$ is the push-forward on operator algebras.
  \item[\textup{(iii)}] \textbf{Composition:}
    $(g \circ f)_* = g_* \circ f_*$.
  \item[\textup{(iv)}] \textbf{Identity:}
    $(\mathrm{id}_X)_* = \mathrm{id}_{\mathfrak{D}(\HC_X)}$.
\end{enumerate}
\end{theorem}

\begin{proof}
\textit{(i)}~Existence of the representation is
Theorem~\ref{thm:5.2}(i)--(ii).

\textit{(ii)}~Given $f : \CM_X \to \CM_Y$ with $f^*g_Y = g_X$
(isometric embedding of constraint manifolds), define
$f_*(\rho) := \Phi_Y(\iota_{Y*} \circ f_* \circ
\iota_X^{-1}(\cdot))(\rho)$ on the level of density operators.
Complete positivity and trace preservation of $f_*$ follow from the
CPTP properties of $\Phi_Y$ (Axiom~\ref{ax:4.3}).

\textit{(iii)}~Functoriality of composition follows from the
chain rule for push-forwards:
$(g \circ f)_* = g_* \circ f_*$ on smooth maps of manifolds, and this
structure is preserved by the CPTP property of the interface maps.

\textit{(iv)}~$(\mathrm{id}_X)_* = \mathrm{id}$ since the identity
map on $\CM_X$ induces the identity push-forward on $\CB(\CR_{\mathrm{op}})$
(Definition~\ref{def:3.17b}).
\end{proof}

% ============================================================
\subsection{Spectral Structure and Superselection}
\label{sec:ch5:spectral}
% ============================================================

\begin{definition}[Superselection Sectors]
\label{def:5.11}
The \emph{superselection sectors} of the cognitive system $X$ are the
irreducible subrepresentations of $\CK(\HC)$ acting on $\HC$:
\[
  \HC = \bigoplus_{\alpha \in \mathcal{I}} \HC^{(\alpha)},
\]
where each $\HC^{(\alpha)}$ is an irreducible $\CK(\HC)$-module, and
$\mathcal{I}$ is a finite or countable index set.
\end{definition}

\begin{proposition}[Decomposition into Superselection Sectors]
\label{prop:5.3}
The cognitive state space $\HC$ decomposes into superselection sectors
with the following properties:
\begin{enumerate}
  \item[\textup{(i)}] Each sector $\HC^{(\alpha)}$ is invariant under
    $\CK(\HC)$: $\CK(\HC) \cdot \HC^{(\alpha)} \subseteq
    \HC^{(\alpha)}$.
  \item[\textup{(ii)}] No observable in $\CA(\HC)$ connects different
    sectors: for $\hatO \in \CA(\HC)$,
    $\langle\psi^{(\alpha)}|\hatO|\psi^{(\beta)}\rangle = 0$ whenever
    $\alpha \neq \beta$ and $\psi^{(\alpha)} \in \HC^{(\alpha)}$,
    $\psi^{(\beta)} \in \HC^{(\beta)}$.
  \item[\textup{(iii)}] The constraint Hamiltonian $\hatH_C$ preserves
    each sector: $\hatH_C \HC^{(\alpha)} \subseteq \HC^{(\alpha)}$.
  \item[\textup{(iv)}] If $\dim\HC < \infty$, the number of sectors
    $|\mathcal{I}|$ is finite and $\sum_\alpha \dim\HC^{(\alpha)} =
    \dim\HC \leq \exp(\kappa_{\max} \cdot \mathrm{vol}(\CM))$.
\end{enumerate}
\end{proposition}

\begin{proof}
(i)~By definition of irreducible subrepresentation.

(ii)~Observables in $\CA(\HC)$ commute with $\hatP_\CM$
(Definition~\ref{def:3.20}).  Since $\CK(\HC) \subseteq \CB(\HC_\CM)$
by construction (the constraint generators map into $\HC_\CM$ via $\Phi$
and the surjectivity condition Axiom~\ref{ax:4.3}(IC2)), the
superselection sectors of $\CK(\HC)$ are respected by $\CA(\HC)$.
Formally: if $P_\alpha$ is the projection onto $\HC^{(\alpha)}$, then
Schur's lemma implies $[A, P_\alpha] = 0$ for all
$A \in \CK(\HC)'$; since $\CA(\HC) \subseteq \CK(\HC)' \cap
\CB(\HC)_{\mathrm{sa}}$ (observables commute with $\hatP_\CM$ and
hence with $\CK(\HC)$ which is built from $\hatP_\CM$-images), the
off-diagonal matrix elements vanish.

(iii)~$\hatH_C = \frac{1}{2}\sum_i \hatC_i^2 \in \CK(\HC)$
(Definition~\ref{def:5.7}), so $\hatH_C$ preserves each
$\CK(\HC)$-invariant subspace.

(iv)~The decomposition is finite in finite dimension by the
Artin--Wedderburn theorem (Proposition~\ref{prop:5.1}(iv)).
The dimension bound is Axiom~\ref{ax:4.4}.
\end{proof}

\begin{lemma}[Superselection and Density Operators]
\label{lem:5.7}
Every density operator $\rho \in \mathfrak{D}(\HC)$ decomposes as
\[
  \rho = \sum_\alpha p_\alpha\, \rho^{(\alpha)},
  \qquad
  p_\alpha := \Tr(P_\alpha\,\rho) \geq 0,
  \quad
  \sum_\alpha p_\alpha = 1,
\]
where $\rho^{(\alpha)} := P_\alpha\,\rho\,P_\alpha / p_\alpha \in
\mathfrak{D}(\HC^{(\alpha)})$ \textup{(}when $p_\alpha > 0$\textup{)}
and $P_\alpha$ is the projection onto the superselection sector
$\HC^{(\alpha)}$.  The off-diagonal blocks
$P_\alpha\,\rho\,P_\beta = 0$ for $\alpha \neq \beta$ are
unobservable.
\end{lemma}

\begin{proof}
For any observable $\hatO \in \CA(\HC)$ and $\alpha \neq \beta$,
$\Tr(P_\alpha\,\rho\,P_\beta\,\hatO) =
\Tr(\rho\,P_\beta\,\hatO\,P_\alpha) = 0$ by
Proposition~\ref{prop:5.3}(ii) (since $\hatO$ is block-diagonal in
the superselection sectors).  The off-diagonal terms
$P_\alpha\,\rho\,P_\beta$ contribute zero to all expectation values
and are hence physically unobservable.  The diagonal blocks give
$\rho = \sum_\alpha P_\alpha\,\rho\,P_\alpha$
(effective superselection), with
$p_\alpha = \Tr(P_\alpha\,\rho)$ and normalised states
$\rho^{(\alpha)}$ in each sector.
\end{proof}

% ============================================================
\subsection{Multi-Agent Tensor Algebra}
\label{sec:ch5:tensor}
% ============================================================

We now develop the algebraic structure of composite cognitive systems.

\begin{definition}[Joint Constraint Operator Algebra]
\label{def:5.12}
For cognitive systems $X_1$ and $X_2$ with constraint operator algebras
$\CK(\HC_1)$ and $\CK(\HC_2)$, the \emph{joint constraint operator
algebra} is
\[
  \CK(\HC_{12}) := \overline{
    \CK(\HC_1) \otimes \CK(\HC_2)
  }^{\|\cdot\|_{\mathrm{op}}}
  \;\subseteq\; \CB(\HC_1 \otimes \HC_2),
\]
where the tensor product is the spatial (minimal) $C^*$-tensor product,
and $\HC_{12} = \HC_1 \otimes \HC_2$ by
Axiom~\textup{\ref{ax:4.6}}(CC1).
\end{definition}

\begin{proposition}[Properties of the Joint Algebra]
\label{prop:5.4}
The joint constraint operator algebra $\CK(\HC_{12})$ satisfies:
\begin{enumerate}
  \item[\textup{(i)}] $\CK(\HC_{12})$ is a unital $C^*$-algebra.
  \item[\textup{(ii)}] The joint constraint generators are
    $\{\hatC_i^{(1)} \otimes \id_2,\; \id_1 \otimes
    \hatC_j^{(2)}\}_{i,j}$.
  \item[\textup{(iii)}] Generators from different subsystems commute:
    $[\hatC_i^{(1)} \otimes \id_2,\;
    \id_1 \otimes \hatC_j^{(2)}] = 0$.
  \item[\textup{(iv)}] The joint constraint Hamiltonian decomposes as
    $\hatH_C^{(12)} = \hatH_C^{(1)} \otimes \id_2
    + \id_1 \otimes \hatH_C^{(2)}$.
\end{enumerate}
\end{proposition}

\begin{proof}
(i)~The minimal $C^*$-tensor product of two $C^*$-algebras is a
$C^*$-algebra.  Unitality: $\id_1 \otimes \id_2 = \id_{12} \in
\CK(\HC_{12})$.

(ii)~By Axiom~\ref{ax:4.6}(CC2), the product constraint operator is
$\hatC_{12} = \hatC_1 \oplus \hatC_2$, which acts on $T\CM_{12} =
T\CM_1 \oplus T\CM_2$.  Applying the interface map $\Phi_{12}$ to the
generators of $\hatC_{12}$ in the product frame
$(e_i^{(1)}, 0) \cup (0, e_j^{(2)})$ yields the stated tensor products
by the product structure of the Stinespring dilation.

(iii)~$(A \otimes \id)(\id \otimes B) = A \otimes B =
(\id \otimes B)(A \otimes \id)$ for all operators $A, B$.

(iv)~$\hatH_C^{(12)} = \frac{1}{2}\sum_i (\hatC_i^{(1)} \otimes
\id_2)^2 + \frac{1}{2}\sum_j (\id_1 \otimes \hatC_j^{(2)})^2
= \bigl(\frac{1}{2}\sum_i (\hatC_i^{(1)})^2\bigr) \otimes \id_2
+ \id_1 \otimes \bigl(\frac{1}{2}\sum_j (\hatC_j^{(2)})^2\bigr)
= \hatH_C^{(1)} \otimes \id_2 + \id_1 \otimes \hatH_C^{(2)}$.
\end{proof}

\begin{lemma}[Entanglement and Constraint Curvature]
\label{lem:5.8}
For a joint system $X_{12}$ with non-zero constraint curvature, the
ground state of $\hatH_C^{(12)}$ can be entangled.  Specifically, if
$\hatH_C^{(1)}$ has ground state $\psi_0^{(1)}$ and
$\hatH_C^{(2)}$ has ground state $\psi_0^{(2)}$, then:
\begin{enumerate}
  \item[\textup{(i)}] If $\hatH_C^{(12)} = \hatH_C^{(1)} \otimes
    \id + \id \otimes \hatH_C^{(2)}$ exactly
    \textup{(}no interaction term\textup{)}, the ground state is the
    product state $\psi_0^{(1)} \otimes \psi_0^{(2)}$ and is
    unentangled.
  \item[\textup{(ii)}] If additional coupling terms arise from the
    sub-multiplicative curvature bound
    Axiom~\textup{\ref{ax:4.6}}(CC3), the effective joint Hamiltonian
    may include interaction terms that entangle the ground state.
\end{enumerate}
\end{lemma}

\begin{proof}
(i)~If $\hatH_C^{(12)}$ has the additive form, then by standard results
on tensor-product Hamiltonians, the ground state is the tensor product
of the individual ground states (assuming non-degeneracy).

(ii)~The sub-multiplicative bound $\kappa_{12}(m_1, m_2) \leq
\kappa_1(m_1) \cdot \kappa_2(m_2)$ (Axiom~\ref{ax:4.6}(CC3)) allows for
non-trivial cross-curvature terms in the joint constraint geometry.  When
$F^{\hatC_{12}}$ has components mixing the $\CM_1$ and $\CM_2$ directions,
the constraint generators of the joint system include interaction terms
$\hatC_i^{(1)} \otimes \hatC_j^{(2)}$ that do not factor, and the
ground state of the resulting Hamiltonian is generically entangled
(by the theory of interacting quantum systems).
\end{proof}

\begin{proposition}[Partial Trace Reduction]
\label{prop:5.5}
For a joint density operator $\rho_{12} \in
\mathfrak{D}(\HC_1 \otimes \HC_2)$, the \emph{reduced state} of
subsystem~$1$ is
\[
  \rho_1 := \Tr_2(\rho_{12}) \in \mathfrak{D}(\HC_1),
\]
where $\Tr_2 : \mathcal{T}_1(\HC_1 \otimes \HC_2) \to
\mathcal{T}_1(\HC_1)$ is the partial trace over $\HC_2$.  This
satisfies:
\begin{enumerate}
  \item[\textup{(i)}] $\rho_1$ is a valid density operator.
  \item[\textup{(ii)}] For any observable $\hatO_1 \in \CA(\HC_1)$,
    $\Tr(\rho_{12}(\hatO_1 \otimes \id_2)) = \Tr(\rho_1 \hatO_1)$.
  \item[\textup{(iii)}] If $\rho_{12} = \rho_1 \otimes \rho_2$
    \textup{(}product state\textup{)}, then
    $\Tr_2(\rho_{12}) = \rho_1$.
\end{enumerate}
\end{proposition}

\begin{proof}
(i)~The partial trace preserves self-adjointness, positivity, and unit
trace:
$\rho_1^\dagger = \Tr_2(\rho_{12}^\dagger) = \Tr_2(\rho_{12}) =
\rho_1$;
$\langle\psi|\rho_1|\psi\rangle = \sum_k \langle\psi \otimes
f_k|\rho_{12}|\psi \otimes f_k\rangle \geq 0$ for any orthonormal
basis $\{f_k\}$ of $\HC_2$;
$\Tr(\rho_1) = \Tr_{12}(\rho_{12}) = 1$.

(ii)~$\Tr(\rho_{12}(\hatO_1 \otimes \id_2))
= \Tr_1(\Tr_2(\rho_{12}(\hatO_1 \otimes \id_2)))
= \Tr_1(\Tr_2(\rho_{12}) \cdot \hatO_1)
= \Tr(\rho_1 \hatO_1)$ by the cyclic property of the trace and the
definition of partial trace.

(iii)~$\Tr_2(\rho_1 \otimes \rho_2) = \rho_1 \cdot \Tr(\rho_2) =
\rho_1$.
\end{proof}

\begin{theorem}[Classical Emergence via Decoherence]
\label{thm:5.4}
In the limit $\hbar_{\mathrm{eff}} \to 0$
\textup{(}equivalently, $\kappa_{\max} \cdot \mathrm{vol}(\CM) \to
\infty$\textup{)}, the algebraic structure of CIIR reduces to a
commutative algebra:
\begin{enumerate}
  \item[\textup{(i)}] $[\hatC_i, \hatC_j] \to 0$ for all $i, j$
    \textup{(}Theorem~\textup{\ref{thm:5.1})}.
  \item[\textup{(ii)}] $\CK(\HC)$ becomes commutative and is
    isomorphic to $C(\Sigma)$ for a compact Hausdorff space $\Sigma$
    \textup{(}the Gel'fand spectrum\textup{)}.
  \item[\textup{(iii)}] The superselection sectors collapse to
    one-dimensional subspaces, and density operators become classical
    probability distributions on $\Sigma$.
  \item[\textup{(iv)}] The Born rule \textup{(}Axiom~\textup{\ref{ax:4.5})}
    reduces to Kolmogorov probability on the measurable space
    $(\Sigma, \mathscr{B}(\Sigma))$.
\end{enumerate}
\end{theorem}

\begin{proof}
(i)~By Theorem~\ref{thm:5.1},
$\|[\hatC_i, \hatC_j]\| = \mathcal{O}(\hbar_{\mathrm{eff}}) \to 0$
as $\hbar_{\mathrm{eff}} \to 0$.

(ii)~A $C^*$-algebra in which all generators commute is commutative.
By the Gel'fand--Naimark theorem, every commutative unital $C^*$-algebra
is isometrically $*$-isomorphic to $C(\Sigma)$ where $\Sigma$ is its
maximal ideal space (a compact Hausdorff space).

(iii)~In a commutative $C^*$-algebra $C(\Sigma)$, every irreducible
representation is one-dimensional (by Schur's lemma for commutative
algebras).  Hence the Hilbert space decomposes as
$\HC = \bigoplus_{\sigma \in \Sigma} \mathbb{C}_\sigma$, and density
operators reduce to probability measures on $\Sigma$.

(iv)~The Born rule $\Tr(E(\Delta)\rho) = \sum_{\sigma \in \Delta}
p_\sigma$ becomes summation of classical probabilities over the outcome
set $\Delta \subseteq \Sigma$.  This is Kolmogorov probability on
$(\Sigma, \mathscr{B}(\Sigma))$.
\end{proof}

% ============================================================
\subsection{Summary of Derived Structures}
\label{sec:ch5:summary}
% ============================================================

\begin{center}
\begin{tabular}{cllp{5.4cm}}
\hline
\textbf{Ref.} & \textbf{Name} & \textbf{Derived from} &
  \textbf{Content} \\
\hline
D~\ref{def:5.1} & Constraint Frame &
  D~\ref{def:3.7} &
  Local frame $\{e_i\}$ on $T\CM$ \\
D~\ref{def:5.2} & Constraint Op.\ Algebra &
  D~\ref{def:3.17b}, D~\ref{def:3.18}, Ax.~\ref{ax:4.3} &
  $\CK(\HC) := \overline{\mathrm{alg}^*\{\Phi(\iota_*(\hatC(e_i)))\}}$ \\
D~\ref{def:5.3} & Constraint Generators &
  D~\ref{def:5.2} &
  $\hatC_i := \Phi(\iota_*(\hatC(e_i)))$ \\
D~\ref{def:5.4} & Holonomy Tensor &
  D~\ref{def:3.9} &
  $\Omega_{ij} := \mathrm{tr}_g(F^{\hatC}(e_i, e_j))$ \\
D~\ref{def:5.5} & Effective $\hbar$ &
  D~\ref{def:3.8}, D~\ref{def:3.6}, Ax.~\ref{ax:4.2} &
  $\hbar_{\mathrm{eff}} := (\kappa_{\max} \cdot
  \mathrm{vol}(\CM))^{-1}$ \\
D~\ref{def:5.6} & Holonomy Operator &
  D~\ref{def:5.4}, D~\ref{def:3.18} &
  $\hat\Omega_{ij} :=
  \Phi(\iota_*(\mathrm{tr}_g(F^{\hatC}(e_i,e_j))))$ \\
D~\ref{def:5.7} & Constraint Hamiltonian &
  D~\ref{def:5.3} &
  $\hatH_C := \frac{1}{2}\sum_i \hatC_i^2$ \\
D~\ref{def:5.8} & Constraint Vacuum &
  D~\ref{def:5.7} &
  $\hatH_C \psi_0 = 0$ \\
D~\ref{def:5.9} & State Functional &
  D~\ref{def:3.14}, D~\ref{def:3.20} &
  $\omega_\rho(A) := \Tr(\rho A)$ \\
D~\ref{def:5.10} & GNS Triple &
  D~\ref{def:5.9} &
  $(\HC_\omega, \pi_\omega, \xi_\omega)$ \\
D~\ref{def:5.11} & Superselection Sectors &
  D~\ref{def:5.2} &
  $\HC = \bigoplus_\alpha \HC^{(\alpha)}$ \\
D~\ref{def:5.12} & Joint Constraint Algebra &
  D~\ref{def:5.2}, Ax.~\ref{ax:4.6} &
  $\CK(\HC_{12}) := \overline{\CK(\HC_1) \otimes \CK(\HC_2)}$ \\
\hline
\end{tabular}
\end{center}

% ============================================================
\subsection{Consistency Check against Chapters 3--4}
\label{sec:ch5:ch34check}
% ============================================================

\noindent\textbf{Verification.}
\begin{enumerate}
  \item \textbf{No new axioms:}
    Every result in this chapter is derived from
    Axioms~\ref{ax:4.1}--\ref{ax:4.6} and
    Definitions~\ref{def:3.1}--\ref{def:3.21}.  No additional postulates
    are introduced.

  \item \textbf{No undefined symbols:}
    Every symbol in Definitions~\ref{def:5.1}--\ref{def:5.12} is either
    a standard mathematical object or is defined in Chapters~3--4.
    The new objects ($\CK(\HC)$, $\hatC_i$, $\Omega_{ij}$,
    $\hbar_{\mathrm{eff}}$, $\hatH_C$, $\psi_0$, $\omega_\rho$,
    $\HC^{(\alpha)}$) are all constructed from the primitives using
    the interface map $\Phi$ and the constraint geometry $(\CM, g,
    \iota, \hatC)$.

  \item \textbf{No circular definitions:}
    The dependency order is:
    $\{e_i\}$ (D~\ref{def:5.1}) $\to$ $\hatC_i$ (D~\ref{def:5.3})
    $\to$ $\CK(\HC)$ (D~\ref{def:5.2}) $\to$ $\Omega_{ij}$
    (D~\ref{def:5.4}) $\to$ $\hbar_{\mathrm{eff}}$ (D~\ref{def:5.5})
    $\to$ $\hat\Omega_{ij}$ (D~\ref{def:5.6}) $\to$ $\hatH_C$
    (D~\ref{def:5.7}) $\to$ $\psi_0$ (D~\ref{def:5.8}) $\to$
    $\omega_\rho$ (D~\ref{def:5.9}) $\to$ GNS (D~\ref{def:5.10})
    $\to$ sectors (D~\ref{def:5.11}) $\to$ joint (D~\ref{def:5.12}).
    This is a directed acyclic graph.

  \item \textbf{No contradiction with Chapters 3--4:}
    The derived structures are consequences of the axioms and primitives.
    The constraint operator algebra $\CK(\HC)$ is a subalgebra of
    $\CB(\HC)$ (Definition~\ref{def:3.12}).  The commutation relations
    are consistent with the Bianchi identity
    (Proposition~\ref{prop:3.2}).  The GNS representation is consistent
    with the observable algebra structure
    (Proposition~\ref{prop:3.6}).  The tensor-product structure is
    consistent with Axiom~\ref{ax:4.6}.

  \item \textbf{Sufficiency Theorem verified:}
    Theorem~\ref{thm:4.3}(i) (algebraic structure) is fully established
    by Definitions~\ref{def:5.2}--\ref{def:5.7} and
    Theorem~\ref{thm:5.1}.  Theorem~\ref{thm:4.3}(ii) (representation
    theorem) is established by Theorem~\ref{thm:5.2}.
    Theorem~\ref{thm:4.3}(iv) (multi-agent structure) is supported by
    Definitions~\ref{def:5.11}--\ref{def:5.12} and
    Propositions~\ref{prop:5.3}--\ref{prop:5.5}.
    Theorem~\ref{thm:4.3}(iii) (dynamics) is deferred to Chapter~6.
\end{enumerate}

% ============================================================
\subsection{Dependency Graph}
\label{sec:ch5:depgraph}
% ============================================================

\begin{center}
\begin{tabular}{lll}
\hline
\textbf{Object} & \textbf{Reference} & \textbf{Depends on} \\
\hline
Constraint frame & Def.~\ref{def:5.1} &
  $(\CM, g)$ (D~\ref{def:3.7}) \\
Constraint generators & Def.~\ref{def:5.3} &
  $\Phi$ (D~\ref{def:3.18}), $\iota_*$ (D~\ref{def:3.17b}),
  $\hatC$ (D~\ref{def:3.7}) \\
Constraint op.\ algebra & Def.~\ref{def:5.2} &
  Constraint generators (D~\ref{def:5.3}) \\
Holonomy tensor & Def.~\ref{def:5.4} &
  $F^{\hatC}$ (D~\ref{def:3.9}), $g$ (D~\ref{def:3.4}) \\
Effective $\hbar$ & Def.~\ref{def:5.5} &
  $\kappa_{\max}$ (D~\ref{def:3.8}), $\mathrm{vol}(\CM)$
  (D~\ref{def:3.6}), Ax.~\ref{ax:4.2} \\
Holonomy operator & Def.~\ref{def:5.6} &
  Holonomy tensor (D~\ref{def:5.4}), $\Phi$ (D~\ref{def:3.18}) \\
Constraint Hamiltonian & Def.~\ref{def:5.7} &
  Constraint generators (D~\ref{def:5.3}) \\
Constraint vacuum & Def.~\ref{def:5.8} &
  $\hatH_C$ (D~\ref{def:5.7}) \\
State functional & Def.~\ref{def:5.9} &
  $\mathfrak{D}(\HC)$ (D~\ref{def:3.14}),
  $\CA(\HC)$ (D~\ref{def:3.20}) \\
GNS triple & Def.~\ref{def:5.10} &
  State functional (D~\ref{def:5.9}) \\
Superselection sectors & Def.~\ref{def:5.11} &
  $\CK(\HC)$ (D~\ref{def:5.2}), $\HC$ (D~\ref{def:3.11}) \\
Joint algebra & Def.~\ref{def:5.12} &
  $\CK(\HC_1)$, $\CK(\HC_2)$ (D~\ref{def:5.2}),
  Ax.~\ref{ax:4.6} \\
\hline
Commutation (T~\ref{thm:5.1}) & Thm.~\ref{thm:5.1} &
  D~\ref{def:5.3}, D~\ref{def:5.5}, D~\ref{def:5.6},
  P~\ref{prop:3.5} \\
CIIR Representation (T~\ref{thm:5.2}) & Thm.~\ref{thm:5.2} &
  D~\ref{def:5.9}, D~\ref{def:5.10},
  Ax.~\ref{ax:4.3}, P~\ref{prop:3.6} \\
Functoriality (T~\ref{thm:5.3}) & Thm.~\ref{thm:5.3} &
  T~\ref{thm:5.2}, D~\ref{def:3.17b}, Ax.~\ref{ax:4.3} \\
Classical emergence (T~\ref{thm:5.4}) & Thm.~\ref{thm:5.4} &
  T~\ref{thm:5.1}, D~\ref{def:5.11}, Ax.~\ref{ax:4.5} \\
\hline
\end{tabular}
\end{center}

\medskip
\noindent
All 12 definitions, 4 theorems, 8 lemmas, 5 propositions, and 1
corollary are complete.  The derived algebraic structure provides the
foundation for the dynamics (Chapter~6) and model class (Chapter~7) of
CIIR.
